\section*{Chapter \ref{ch:iv1:timed-test-6} --- Timed Test 6}
\markboth{\protect\omTocMark{}{Answers: Timed Test 6}}{\protect\omTocMark{}{Answers: Timed Test 6}}

\begin{solution}{iq:iv1:timed-test-6:1}
The rook attacks with probability $\frac29$, and the king with probability $\frac{5}{48}$.

A rook attacks the same 14 squares wherever it stands: 7 in its row and 7 in its column. The other piece is on one of the 63 remaining squares, so the chance is $\frac{14}{63} = \frac29$.

A king's reach depends on where it stands\index{conditioning on position}. There are 4 corners with 3 neighbours, 24 edge squares with 5, and 36 inner squares with 8. The neighbouring ordered pairs total $12 + 120 + 288 = 420$, out of $64 \times 63 = 4032$. So the chance is $\frac{5}{48}$. An enumeration confirms both.

\iqwrong{Using 8 neighbours everywhere, which gives $\frac{8}{63}$. It ignores that the board's edge cuts off about a fifth of those neighbours, since 28 of the 64 squares have fewer than 8.}
\end{solution}

\begin{solution}{iq:iv1:timed-test-6:2}
The first probability is $\frac12$, and the second is $\frac38$.

One team wins its 3 games with chance $\frac18$. Two teams cannot both do so, because they meet. So these events are disjoint\index{inclusion--exclusion}, and the sum is $4 \times \frac18$.

For the second, choose the winner and the loser, in $4 \times 3$ ordered ways. Together, these two teams play 5 distinct games, as their match counts once. So each ordered choice has chance $2^{-5}$. These events are disjoint too, which gives $\frac{12}{32} = \frac38$.

An enumeration over all $2^6$ outcomes confirms both. With $n$ teams, the first chance is $n2^{-(n - 1)}$, which goes to 0 quickly, as it should.

\iqwrong{Using $1 - (\frac78)^4$, which treats the four events as independent. They exclude each other, so the plain sum is exact.}
\end{solution}

\begin{solution}{iq:iv1:timed-test-6:3}
The expected position is $\frac{53}{5} = 10.6$, and the probability is $\frac{1}{17}$.

The 4 aces split the 48 other cards into 5 gaps. By symmetry, all gaps have the same expected size\index{symmetry}, $\frac{48}{5} = 9.6$. The first ace comes after the first gap, so its position has mean $10.6$.

Whatever the first card is, 3 of the remaining 51 cards share its rank, and $\frac{3}{51} = \frac{1}{17}$.

A simulation agrees with both. As a limit, with 48 aces and 4 other cards, the same rule gives $\frac{53}{49}$, just above 1, as it should be.

\iqwrong{Answering 13, the deck size over the number of aces. This is the spacing between aces if the deck is a circle, and it ignores the gap after the last ace, which also takes its share.}
\end{solution}

\begin{solution}{iq:iv1:timed-test-6:4}
The call is worth $\frac{15}{4} = 3.75$, replicated with $\frac34$ of a share and $-33.75$ in cash. A real probability of 90\% does not change the value.

The call pays 15 or 0. Holding $\Delta$ shares and cash $c$, we need $65\Delta + c = 15$ and $45\Delta + c = 0$\index{replication}. So $\Delta = \frac{15}{20} = \frac34$ and $c = -33.75$. The cost is $\frac34 \times 50 - 33.75 = 3.75$.

As a second route, the risk-neutral probability $p$ makes the stock fair: $65p + 45(1 - p) = 50$, so $p = \frac14$, and the value is $\frac14 \times 15$.

Any other price gives a free profit against the replicating portfolio, whatever the real odds. The real expected payoff of $13.5$ is not a price.

\iqwrong{Pricing at the real expected payoff. This is the usual error, and the replication shows why it fails.}
\end{solution}

\begin{solution}{iq:iv1:timed-test-6:5}
The expected number of neighbouring same-colour pairs is 25, and the expected number of pairs in a hand is $\frac{10}{17}$.

Each of the 51 neighbouring positions has a same-colour pair with chance $\frac{25}{51}$. Indicators add\index{indicator variables} with no need for independence, so the mean is exactly 25.

In a hand, each of the $\binom52 = 10$ pairs of cards has the same rank with chance $\frac{3}{51}$. So the mean is $\frac{30}{51} = \frac{10}{17} \approx 0.59$.

A simulation returns 25.0 and $0.588$. Note that a hand with three of a kind counts as 3 pairs here, and the indicators count it correctly.

\iqwrong{Using $\frac12$ for the neighbouring pair, which gives $25.5$. Having seen a red, there is one fewer red left.}
\end{solution}

\begin{solution}{iq:iv1:timed-test-6:6}
The expected number of lines is 2, and the probability is $\frac{37}{64}$.

There are 8 lines. Each is single-coloured with chance $2 \times 2^{-3} = \frac14$. Indicators give\index{indicator variables} $8 \times \frac14 = 2$, even though the lines overlap.

The three rows share no squares, so they are independent. No row is single-coloured with chance $(\frac34)^3 = \frac{27}{64}$, and the answer is the complement.

An enumeration over the 512 grids confirms both. Only 32 of the 512 grids have no such line at all. Dependence moves the distribution of the count, but not its mean.

\iqwrong{Taking the chance of at least one line among all 8 as $1 - (\frac34)^8$. The lines share squares, so they are not independent. Only the rows are.}
\end{solution}

\begin{solution}{iq:iv1:timed-test-6:7}
The expected number is about $3.49$, and the probability is about $0.0163$.

A given bin is empty with chance $0.9^{10}$. So the mean is $10 \times 0.349 = 3.49$\index{indicator variables}, a little below the large-$n$ value of $10/e = 3.68$.

With exactly one empty bin, one bin has two balls and the other eight have one each. Choose the empty bin in 10 ways. Choose the two balls that share, in $\binom{10}{2} = 45$ ways. Assign the 9 groups to 9 bins, in $9!$ ways. So
\[
P = \frac{10 \times 45 \times 362\,880}{10^{10}} = 0.0163.
\]

A simulation confirms both figures. With $3.49$ empty bins on average, having only one is rare, and that is what this small value says.

\iqwrong{Multiplying $0.349$ by the chance that the rest are full, as if the bins were independent. Filling nine bins with ten balls uses a precise count instead.}
\end{solution}

\begin{solution}{iq:iv1:timed-test-6:8}
The probability is $\frac{244}{495} \approx 0.493$.

The first roll wins at once with chance $\frac{6 + 2}{36} = \frac29$.

For a point with $w$ ways of rolling it, only the point or a 7 matters on later rolls. So the point comes first with chance $\frac{w}{w + 6}$\index{conditioning on the relevant event}. The points 4 and 10 have $w = 3$, 5 and 9 have $w = 4$, and 6 and 8 have $w = 5$. Thus
\[
P = \frac29 + 2\left(\frac{3}{36} \cdot \frac39 + \frac{4}{36} \cdot \frac{4}{10} + \frac{5}{36} \cdot \frac{5}{11}\right).
\]
The bracket is $\frac{1}{36} + \frac{2}{45} + \frac{25}{396}$. Adding all this up gives $\frac{244}{495}$. A simulation gives $0.493$. So the house edge is $1.4\%$ on an even-money bet.

\iqwrong{Using $\frac{w}{36}$ for the point's chance on the later rolls. Rolls that show neither the point nor a 7 do not count.}
\end{solution}

\begin{solution}{iq:iv1:timed-test-6:9}
The volatility is about $10.6\%$, and the limit is about $8.9\%$.

For $n$ equally weighted assets with variance $\sigma^2$ and correlation $\rho$\index{diversification},
\[
\sigma_p^2 = \sigma^2\left(\frac{1}{n} + \frac{n - 1}{n}\rho\right).
\]
With $n = 10$, the bracket is $0.1 + 0.9 \times 0.2 = 0.28$. So the variance is $0.04 \times 0.28 = 0.0112$, and the volatility is $\sqrt{0.0112} = 10.6\%$.

As $n$ grows, the bracket tends to $\rho$, which gives $\sqrt{0.008} = 8.9\%$. This is the risk that diversification cannot remove.

A check at the limits: $\rho = 1$ gives $20\%$ whatever $n$ is, and $\rho = 0$ gives $20\%/\sqrt{10} = 6.3\%$.

\iqwrong{Dividing by $\sqrt{10}$ and stopping at $6.3\%$, which needs independence. Most of the risk left comes from the common correlation.}
\end{solution}

\begin{solution}{iq:iv1:timed-test-6:10}
The expected length is $\frac12 + x - x^2$, which is $\frac34$ at $x = \frac12$ and $\frac{11}{16}$ at $x = \frac14$.

If the cut $U$ falls left of $x$, then the piece is $(U, 1)$, with length $1 - U$. If it falls right of $x$, the piece is $(0, U)$, with length $U$. So
\[
\mathbb{E} = \int_0^x(1 - u)\,du + \int_x^1u\,du = x - \frac{x^2}{2} + \frac12 - \frac{x^2}{2}.
\]

Limit check: at $x = 0$, the piece always starts at 0, and its length $U$ has mean $\frac12$, which matches. A simulation gives the two values.

Note that the piece containing a fixed point is longer than an average piece, $\frac12$. This is length-biased sampling\index{length-biased sampling}, which also underlies the inspection paradox.

\iqwrong{Answering $\frac12$ for all $x$. The point is more likely to fall in the longer of the two pieces.}
\end{solution}

\begin{solution}{iq:iv1:timed-test-6:11}
The length is uniform on $1, \dots, 10$, with mean $5.5$, and the probability is $\frac12$.

Follow element 1 under the permutation. It returns to itself on the first step with chance $\frac{1}{10}$. If it does not, it does so on the next step with chance $\frac19$, and so on. The chance of a length $k$ is
\[
\frac{9}{10} \cdot \frac89 \cdots \frac{11 - k}{12 - k} \cdot \frac{1}{11 - k} = \frac{1}{10}.
\]
So it is uniform\index{cycle structure}. Given a length $L$, element 2 is among the other $L - 1$ members of the cycle with chance $\frac{L - 1}{9}$. Its mean over $L$ is $\frac{4.5}{9} = \frac12$.

A simulation agrees. A limit check: with $n = 2$, the two elements are together only under the swap, which has chance $\frac12$. The value does not depend on $n$.

\iqwrong{Assuming that long cycles are rarer for a given element. A permutation has more short cycles than long ones, but a long cycle holds many more elements, and the two effects balance exactly.}
\end{solution}

\begin{solution}{iq:iv1:timed-test-6:12}
Roll twice, keep 21 of the 25 outcomes, and repeat otherwise. That costs $\frac{50}{21} \approx 2.38$ rolls per output, against a bound of $\log_57 \approx 1.21$.

Two rolls give 25 equally likely pairs. Give each of the 7 outputs three pairs, using 21 pairs, and start over on the other 4\index{rejection sampling}. Each output then has the same chance, so the result is fair.

Each attempt succeeds with chance $\frac{21}{25}$, so the number of attempts is geometric with mean $\frac{25}{21}$. Each attempt costs 2 rolls. A simulation returns $2.38$.

Each output carries $\log_27$ bits, and each roll $\log_25$ bits, so no scheme can use fewer than $\log_57$ rolls per output on average.

\iqwrong{Summing two rolls and reducing modulo 7. The sums are not equally likely, so the output is not fair, however it looks.}

\iqfollow{Can the 4 rejected outcomes be reused?}{Yes. They carry information: a rejected pair is a uniform draw from 4, which can seed the next attempt. This brings the cost closer to the bound.}
\end{solution}
